% concatenated from squarefreness-polynomials.tex
\documentclass[11pt]{amsart}
\pdfoutput=1
\usepackage{verbatim}

\usepackage[dvipsnames, table, xcdraw]{xcolor}
\usepackage{upgreek}
\usepackage{amssymb,amsmath,mathtools,amsfonts,amsthm}
\usepackage[mathscr]{eucal}
\usepackage{fullpage}
\usepackage{color}
\usepackage{caption}
\usepackage{subcaption}
\usepackage{enumitem}
\setlist[itemize]{leftmargin=30pt, itemsep=2pt}
\setlist[enumerate]{leftmargin=30pt, itemsep=2pt}
\usepackage{soul}
\usepackage{setspace} \setstretch{1.2} 
\usepackage[numbers]{natbib}
\setlength{\bibsep}{0.0pt}
\usepackage{indentfirst} 

\usepackage[colorlinks,pagebackref]{hyperref}
\renewcommand\backrefxxx[3]{\hyperlink{page.#1}{$\uparrow$#1}}
\definecolor{mylinkcolor}{RGB}{0,0,255}
\definecolor{mycitecolor}{RGB}{169,169,169}
\definecolor{myurlcolor}{RGB}{255,20,147}
\definecolor{mybiburlcolor}{RGB}{80,80,80}
\hypersetup{
  colorlinks=true,
  urlcolor=myurlcolor,
  citecolor=mycitecolor,
  linkcolor=mylinkcolor,
  linktoc=page,
  breaklinks=true
}
\usepackage{tikz-cd,tikz} 
\usepackage{cleveref} 

% --------------------------------------------------
% algorithms
% --------------------------------------------------
\usepackage{algorithmic}
\usepackage[ruled]{algorithm}
\makeatletter 
\renewcommand\thealgorithm{\thesection.\arabic{algorithm}} 
\@addtoreset{algorithm}{section} 
\makeatother

\numberwithin{equation}{section} 
\newtheorem{thmx}{Theorem} \renewcommand{\thethmx}{\Alph{thmx}}
\newtheorem{corox}{Corollary} \renewcommand{\thecorox}{\Alph{corox}}
\newtheorem*{maintheorem}{Theorem}
\newtheorem{theorem}[algorithm]{Theorem}
\newtheorem{lemma}[algorithm]{Lemma}
\newtheorem{coro}[algorithm]{Corollary}
\newtheorem{conjecture}[algorithm]{Conjecture}
\newtheorem{proposition}[algorithm]{Proposition}

\theoremstyle{definition} 
\newtheorem{defn}[algorithm]{Definition}
\newtheorem{question}[algorithm]{Question}
\newtheorem{remark}[algorithm]{Remark}
\newtheorem{problem}[algorithm]{Problem}
\newtheorem{example}[algorithm]{Example}

\crefname{equation}{equation}{equations}
\Crefname{equation}{Equation}{Equations}
\crefname{conjecture}{conjecture}{conjectures}
\Crefname{conjecture}{Conjecture}{Conjectures}

\newcommand{\RR}{\mathbb{R}} 
\newcommand{\ZZ}{\mathbb{Z}}
\newcommand{\QQ}{\mathbb{Q}} 
\newcommand{\CC}{\mathbb{C}} 
\newcommand{\PP}{\mathbb{P}} 
\renewcommand{\AA}{\mathbb{A}} 
\newcommand{\cE}{\mathcal{E}}
\newcommand{\cF}{\mathcal{F}} 
\newcommand{\cM}{\mathcal{M}} 
\newcommand{\cB}{\mathcal{B}} 
\newcommand{\cA}{\mathcal{A}}
\newcommand{\GG}{\mathbb{G}} 
\newcommand{\cO}{\mathcal{O}} 
\newcommand{\gor}{\mathsf{gor}} 
\newcommand{\cS}{\mathcal{S}} 
\newcommand{\cN}{\mathcal{N}} 
\newcommand{\cH}{\mathcal{H}} 
\newcommand{\ove}{\overline{e}} 
\newcommand{\ovf}{\overline{f}} 
\newcommand{\ovd}{\overline{d}}
\newcommand{\liesl}{\mathfrak{sl}}
\newcommand{\liegl}{\mathfrak{gl}}
\newcommand{\cI}{\mathcal{I}} 

\newcommand{\overbar}[1]{\mkern 1.5mu\overline{\mkern-1.5mu#1\mkern-1.5mu}\mkern 1.5mu}

\DeclareMathOperator{\Log}{Log} 
\DeclareMathOperator{\GO}{GO} 
\DeclareMathOperator{\disc}{disc} 
\DeclareMathOperator{\Aut}{Aut}
\DeclareMathOperator{\codim}{codim} 
\DeclareMathOperator{\rank}{rk} 
\DeclareMathOperator{\Spec}{Spec}
\DeclareMathOperator{\Proj}{Proj} 
\DeclareMathOperator{\Sym}{Sym} 
\DeclareMathOperator{\Tsch}{Tsch} 
\DeclareMathOperator{\Hom}{Hom} 
\DeclareMathOperator{\Bun}{Bun} 
\DeclareMathOperator{\polytope}{\mathcal{P}} 
\DeclareMathOperator{\qpolytope}{\mathcal{Q}} 
\DeclareMathOperator{\ev}{ev} 
\DeclareMathOperator{\End}{End}
\DeclareMathOperator{\diag}{diag}
\DeclareMathOperator{\pr}{pr}
\DeclareMathOperator{\GL}{GL}
\DeclareMathOperator{\SL}{SL}
\DeclareMathOperator{\PGL}{PGL}
\renewcommand{\subset}{\subseteq}
\newcommand{\mymat}[4]{\ensuremath{\left( \begin{smallmatrix} #1 & #2 \\ #3 & #4 \end{smallmatrix} \right)}}
\newcommand{\marksing}{\mathcal{U}}

\newcommand{\I}{I}
\newcommand{\J}{J}
\newcommand{\K}{K}
\renewcommand{\L}{L}
\newcommand{\fatemeh}[1]{\textcolor{cyan}{Fatemeh: #1}}

%%%%%%%%%%%%%%%%%%%%%%%%%%%%%%%%%%%%%%%%%%%%%%%%%%%%%%%%%%%%%%%%%%%%%%%%%%%%%%%%
%%%%%%%%%%%%%%%%%%%%%%%%%%%%%%%%%%%%%%%%%%%%%%%%%%%%%%%%%%%%%%%%%%%%%%%%%%%%%%%%
%%%%%%%%%%%%%%%%%%%%%%%%%%%%%%%%%%%%%%%%%%%%%%%%%%%%%%%%%%%%%%%%%%%%%%%%%%%%%%%%
\title[On the Squarefree Values of Degree-$\boldsymbol{2q}$ Polynomials]{On the Squarefree Values of Degree-$\boldsymbol{2q}$ Polynomials}

\author{Sergio Ricardo Zapata Ceballos}
\address{Dept. of Mathematics and Statistics, Youngstown State University, Youngstown, OH~44555}
\email{srzapataceballos@ysu.edu}

\author{Fatemeh Jalalvand}
\address{Dept. of Mathematics, University of Calgary, Calgary, AB~T2N 1N4}
\email{fatemeh.jalalvand@ucalgary.ca}


\begin{document}
\maketitle



\begin{abstract}
    In this work, we show that if $h(x)$ is an irreducible monic integer polynomial of degree $2q$ (with $q$ prime), whose defining field extension of $\QQ$ contains a Galois extension of degree $q$, then there is a positive density of integers $n$ such that $h(n)$ is squarefree; in particular, $h(n)$ is squarefree for infinitely many integers $n$. As an application, we prove that the family of exceptional cubic fields contains an infinite subfamily whose unit shapes converge to the hexagonal lattice. To the best of our knowledge, this is the first example of a family of non-Galois totally real cubic fields whose unit shapes converge to the hexagonal lattice.
\end{abstract}

\section{Introduction}

Let $h(x)\in\mathbb{Z}[x]$, and let $S$ denote the set of positive integers $n$ for which $h(n)$ is squarefree; that is, $p^2\nmid h(n)$ for every prime $p$. Under suitable mild hypotheses on $h$, it is conjectured that $S$ is infinite, or, more strongly, that $S$ has positive density. This conjecture has proved difficult in general and has only been verified  for $\deg h = 1,2,3$. The cases $\deg h=1$ and $\deg h=2$ can be handled by elementary sieve methods, whereas the case $\deg h=3$ requires a more complicated  argument~\cite{erdos,hooley}. The cases $\deg h\geq4$ remain open in general. Nevertheless, Granville~\cite{granville} showed that the conjecture follows from the \textit{abc} conjecture.

In this paper, we establish the existence of a positive density of squarefree values for a family of polynomials of degree $2q$, where $q$ is prime. Our main result is the following.

\begin{theorem}
\label{thm: main}
    Let $L/\mathbb{Q}$ be a number field of degree $2q$, where $q$ is prime, and suppose that $L=\mathbb{Q}(\alpha)$, where $\alpha$ has irreducible monic polynomial $h(x)\in\mathbb{Z}[x]$. If $L$ contains a Galois subextension $K/\mathbb{Q}$ of degree $q$ and $\gcd\{h(n) \vcentcolon n \in \ZZ\}$ is squarefree, then there is a positive density of integers $n$ such that $h(n)$ is squarefree. 
\end{theorem}

The hypothesis on the subfield $K$ is particularly useful because it gives $[L:K]=2$.
Consequently, the relative norm
\[
    N_{L/K}(x - \alpha)
\]
is a quadratic polynomial over $K$. Moreover, for every integer $n$,
\[
   N_{K/\mathbb{Q}}\bigl(N_{L/K}(n- \alpha)\bigr)=h(n).
\]
Thus, the problem of finding squarefree values of $h$ can be approached by first studying the quadratic polynomial $N_{L/K}(x-\alpha)$ over the ring of integers of $K$. We apply an elementary sieve to the prime ideals dividing the ideal
\[
   \bigl(N_{L/K}(n-\alpha)\bigr).
\]

We then transfer these estimates to the integer $h(n)$ via the norm map. The prime ideal factorization of $\bigl(N_{L/K}(n-\alpha)\bigr)$ determines the prime ideal factorization of $(h(n))$ after taking norms to $\mathbb{Q}$. To ensure that squarefreeness is preserved under this passage, two additional properties must be established. First, the prime ideals dividing $\bigl(N_{L/K}(n-\alpha)\bigr)$ must have residue degree one over $\mathbb{Q}$; see Proposition~\ref{prop: residue_degree}. Second, no two distinct Galois conjugates of a given prime ideal may simultaneously divide $\bigl(N_{L/K}(n-\alpha)\bigr)$; see Proposition~\ref{prop: conjugatenotdividing}. These properties allow us to control the contribution of each rational prime to $h(n)$ and, ultimately, to deduce the positive-density statement in the theorem.

As an immediate corollary, irreducible quartic monic integer polynomials whose Galois group is neither $S_4$ nor $A_4$ (Corollary~\ref{coro:quartics}), take infinitely many squarefree values.

In the remainder of the paper, we apply Theorem~\ref{thm: main} to construct a family of cubic fields for which the shapes of their unit lattices converge to the hexagonal lattice (Theorem~\ref{thm: shapes-convergence}). Briefly, the \emph{unit lattice} of a number field is the image of the unit group of its ring of integers under the \emph{logarithmic embedding} into $\mathbb{R}^n$, for a suitable $n$. The \emph{shape} of this lattice is its equivalence class under scaling, rotation, and reflection. The same notion may also be defined for arbitrary orders. In recent years, there has been growing interest in the shapes of unit lattices~\cite{D4unit,harron,cusick,dang,david, CorsoRodriguezHertz2025, cruz2025classificationtotallyrealnumber}. 

Nevertheless, very little is known about the shapes of maximal orders. Some previous approaches construct orders that are expected to be maximal, but establishing maximality requires a difficult squarefree sieve. In the cubic case, the unit shape of a totally real Galois cubic field is known to be the hexagonal lattice, whereas the corresponding shapes for non-Galois totally real cubic fields remain largely unexplored. Our approach begins with a pair of units that Ennola conjectured (Conjecture~\ref{conj: enola}) to form a fundamental system of units for a family of non-Galois totally real cubic fields. Theorem~\ref{thm: main} allows us to verify this conjecture for a positive proportion of the fields in this family. We then prove that the shapes of the corresponding unit lattices converge to the hexagonal lattice.


\section{Proof of Theorem~\ref{thm: main}}

First, we introduce some notation that will be used throughout this section. Let $q$ be a prime number, and let $h(x)\in\ZZ[x]$ be an irreducible monic polynomial of degree $2q$. Let $\alpha$ be a root of $h(x)$ and set $L \coloneqq \QQ(\alpha)$. Suppose $L$ contains a Galois subextension $K/\QQ$ of degree $q$, and let $\beta$ be the Galois conjugate of $\alpha$ over $K$. We let $\cO_K$ denote the ring of integers of $K$. Since \(h\) is monic, \(\alpha\) is an algebraic integer, and so is its conjugate \(\beta\). Hence, we define
\[
    s \coloneqq \alpha  + \beta \in\mathcal{O}_K,\qquad t \coloneqq \alpha\beta\in\cO_K,
\]
so that \(m(x) := x^2 - sx + t \in\cO_K[x]\). Observe that
\[
    N_{L/K}(n - \alpha)=(n - \alpha)(n  - \beta) = m(n).
\]

\begin{proposition}
\label{prop: residue_degree}
   Let $\mathfrak p$ be a prime ideal of $\mathcal O_K$, and write $\mathfrak p\cap\ZZ=p\ZZ$. If, for some $n\in\ZZ$,
$
\mathfrak p \mid (N_{L/K}(n - \alpha))
$
and $p\nmid Disc(h)$, then $f_{\mathfrak p/p}=1$.
\end{proposition}
\begin{proof}
    By assumption, we have 
    $
        m(n)\equiv 0 \pmod{\mathfrak p}.
    $
    Since \(K/\QQ\) is Galois of degree \( q \), the residue degree
    \(f_{\mathfrak p/p}\) is either \( 1 \) or \( q \). For the sake of contradiction, suppose
    that \(f_{\mathfrak p/p} = q\). Then \( p \) is inert in \(K\), so
    \(\mathfrak p\) is the unique prime of \(\mathcal{O}_K\) above \(p\). In particular,
    every automorphism of \(K/\QQ\) fixes \(\mathfrak p\).
    Let \(\sigma\) generate \(\operatorname{Gal}(K/\QQ)\). For
    \(m(x)=x^2-sx+t\), write
    \[
        m^\sigma(x)=x^2-\sigma(s)x+\sigma(t),
    \]
    and similarly for \(m^{\sigma^i}(x)\) with $ i \in \{1,2,\dots, q\}$. Since
    \(\sigma(\mathfrak p)=\mathfrak p\), applying \(\sigma^i\) to
    \(m(n)\equiv0\pmod{\mathfrak p}\) yields
    \[
        m^{\sigma^i}(n)\equiv0\pmod{\mathfrak p},\qquad \forall{i \in \{1,2,\dots, q\}}
    \]
    Set
    \[
        F(x)=\prod_{i =1}^qm^{\sigma^i}(x).
    \]
    The automorphism \(\sigma\) cyclically permutes the $q$ factors, hence
    \(F(x)\in\QQ[x]\). Since \(F\) is monic of degree \(2q\) and \(F(\alpha)=0\),
    the irreducibility of \(h\) gives
    \begin{equation}
    \label{eq:h-factorization}
        h(x) = F(x) = \prod_{i =1}^qm^{\sigma^i}(x).
    \end{equation}
    Reducing modulo \(\mathfrak p\), the element \(\bar n\) is a root of each of
     \(\overline{m^\sigma}(x),\dots, \overline{m^{\sigma^q}}(x) \). Therefore
    \[
        (x-\bar n)^q\mid \bar h(x)
    \]
    in \((\mathcal{O}_K/\mathfrak p)[x]\). In particular,
    \[
        \bar h(\bar n)=0,\qquad \bar h'(\bar n)=0.
    \]
    Since \(h,h'\in\ZZ[x]\) and \(n\in\ZZ\), this implies
    \[
        h(n)\equiv0\pmod p,\qquad h'(n)\equiv0\pmod p.
    \]
    Thus \(\bar h\) and \(\bar h'\) have a common root over \(\mathbb F_p\), so
    \(\bar h\) is not squarefree modulo \(p\). Equivalently,
        $p\mid\operatorname{Disc}(h)$,
    contrary to the hypothesis. Hence \(f_{\mathfrak p/p}\neq q\), and since $K$ is a prime degree Galois subfield, the
    only possible residue degrees are \(1\) and \(q\), we conclude that
    \(f_{\mathfrak p/p}=1\).
\end{proof}

\begin{proposition}
\label{prop: conjugatenotdividing}
      Let \(\mathfrak{p}\) be a prime in $\cO_K$ and $n \in \ZZ$, with $p=\mathfrak{p}\cap \mathbb{Z}$  such that $\mathfrak{p} \mid (N_{L/K}(n - \alpha))$ and $p\nmid Disc(h)$. If $\bar{\mathfrak{p}}$ is a conjugate of $\mathfrak p$ and $\bar{\mathfrak{p}} \neq \mathfrak{p} $, then  $\bar{\mathfrak{p}} \nmid (N_{L/K}(n - \alpha))$.
\end{proposition}
\begin{proof}
    Assume, for the sake of contradiction, that
    \[
        \mathfrak{p}\mid (m(n)) \qquad\text{and}\qquad \bar{\mathfrak{p}}\mid (m(n)).
    \]
    Let \(\sigma\) generate \(\operatorname{Gal}(K/\QQ)\). Since $\bar{\mathfrak{p}}\neq\mathfrak{p}$, there exists some
    $j\in\{1,\dots,q-1\}$ such that $\bar{\mathfrak{p}}=\sigma^j(\mathfrak{p})$. It follows that both $m(n)$ and $m^{\sigma^{-j}}(n)$ belong to $\mathfrak{p}$. Using the factorization of $h(x)$ given in Equation~\ref{eq:h-factorization} and reducing modulo $\mathfrak{p}$, we obtain
    \[
        (x-\bar{n})^2\mid\overline{h}(x).
    \]
    Hence, $\overline{h}(x)$ has a repeated root modulo $p$, implying that $p\mid\operatorname{Disc}(h)$, contrary to our assumption. Therefore,
    \[
        \bar{\mathfrak{p}}\nmid\bigl(N_{L/K}(n-\alpha)\bigr).
    \]
\end{proof}

Let $\mathcal{I}$ be the set of prime divisors of $\operatorname{Disc}(h)$ and $\mathcal{I}_K$ be the set of primes of $\cO_K$ above the primes in $\mathcal{I}$.

\begin{proposition}
\label{prop: sieve}
Suppose that $\gcd\{h(n):n\in\mathbb{Z}\}$ is squarefree. Then, for each $p\in\mathcal{I}$, there exists an integer $r_p$ such that $h(r_p)\not\equiv 0\pmod{p^2}$. Moreover, there is a positive proportion of integers $n$ satisfying both
    \[
        n\equiv r_p\pmod{p^2\mathcal{O}_K}
        \qquad\text{for all }p\in\mathcal{I},
    \]
and
    \[
        \bigl(N_{L/K}(n-\alpha)\bigr)
    \]
is squarefree away from $\mathcal{I}_K$.
\end{proposition}

\begin{proof}
    For each prime ideal $\mathfrak{p}$ of $\mathcal{O}_K$, define
    \[
        \rho(\mathfrak{p}^2) \coloneqq \#\left\{\overline{n}\in\mathcal{O}_K/\mathfrak{p}^2
        \vcentcolon n\in\mathbb{Z} \text{ and } m(n)\equiv 0\pmod{\mathfrak{p}^2}\right\}.
    \]
    We first estimate $\rho(\mathfrak{p}^2)$. Suppose that
    $\mathfrak{p}\nmid 2\operatorname{Disc}(m)$. Then $m(x)$ has either no roots or two distinct roots modulo $\mathfrak{p}$. By Hensel's lemma, each root modulo $\mathfrak{p}$ lifts uniquely to a root modulo $\mathfrak{p}^2$. Hence
    \[
        \rho(\mathfrak{p}^2)\leq 2 
    \]
    for every $\mathfrak{p}\nmid 2\operatorname{Disc}(m)$. On the other hand, if
    $\mathfrak{p}\mid 2\operatorname{Disc}(m)$, then
    \[
        \rho(\mathfrak{p}^2)<p^2,
    \]
    where $p$ is the rational prime below $\mathfrak{p}$. Indeed, otherwise $m(n)\equiv0\pmod{\mathfrak{p}^2}$ for every $n\in\mathbb{Z}$, and hence
    \[
        h(n) = N_{K/\mathbb{Q}}(m(n)) \equiv0\pmod{p^2}
    \]
    for every $n\in\mathbb{Z}$, contradicting the assumption that $\gcd\{h(n):n\in\mathbb{Z}\}$ is squarefree.
    
    Fix $X\geq1$, and define
    \[
        \mathcal{P}_X \coloneqq \left\{\mathfrak{p}\notin\mathcal{I}_K \vcentcolon N_{K/\mathbb{Q}}(\mathfrak{p}^2)\leq X \right\}.
    \]
    Set
    \[ 
        P_X\coloneqq \prod_{\mathfrak{p}\in\mathcal{P}_X}\mathfrak{p}^2
        \qquad\text{and}\qquad M\coloneqq \prod_{p\in\mathcal{I}}p^2\mathcal{O}_K.
    \]
    Finally, let
    \[
        Z_X \coloneqq \left\{\overline{n}\in\mathcal{O}_K/MP_X \vcentcolon n\in\mathbb{Z}\right\}.
    \]
    
    Since $\gcd\{h(n):n\in\mathbb{Z}\}$ is squarefree, for each $p\in\mathcal{I}$ there exists an integer $r_p$ such that
    \[
        h(r_p)\not\equiv0\pmod{p^2}.
    \]
    
    For each $\mathfrak{p}\in\mathcal{P}_X$, define
    \[
        N_{\mathfrak{p}} \coloneqq \left\{\overline{n}\in\mathcal{O}_K/\mathfrak{p}^2 \vcentcolon n\in\mathbb{Z} \text{ and } m(n)\not\equiv0\pmod{\mathfrak{p}^2}\right\}.
    \]
    We then define
    \[
        N_X \coloneqq \left\{\overline{n}\in Z_X \vcentcolon
        \begin{array}{l}
            n\equiv r_p\pmod{p^2\mathcal{O}_K}
            \text{ for all }p\in\mathcal{I},\\[2mm]
            m(n)\not\equiv0\pmod{\mathfrak{p}^2}
            \text{ for all }\mathfrak{p}\in\mathcal{P}_X
        \end{array}
        \right\}.
    \]
    
    By the Chinese remainder theorem, the natural ring homomorphism
    \[
        \phi: \mathcal{O}_K/MP_X \longrightarrow \prod_{p\in\mathcal{I}} \mathcal{O}_K/p^2\mathcal{O}_K \times \prod_{\mathfrak{p}\in\mathcal{P}_X} \mathcal{O}_K/\mathfrak{p}^2
    \]
    is an isomorphism. By Proposition~\ref{prop: conjugatenotdividing}, the prime ideals occurring in the product are pairwise non-conjugate, and hence
    \[
        \phi(N_X) = \prod_{p\in\mathcal{I}}\{\overline{r_p}\} \times \prod_{\mathfrak{p}\in\mathcal{P}_X} N_{\mathfrak{p}}.
    \]
    Therefore,
    \[
        \#N_X = \prod_{\mathfrak{p}\in\mathcal{P}_X} \left(p^2-\rho(\mathfrak{p}^2)\right) = \left(\prod_{\mathfrak{p}\in\mathcal{P}_X}p^2\right) \prod_{\mathfrak{p}\in\mathcal{P}_X} \left(1-\frac{\rho(\mathfrak{p}^2)}{p^2}\right),
    \]
    where $p$ denotes the rational prime below $\mathfrak{p}$. Since every $\mathfrak{p}\in\mathcal{P}_X$ is unramified, we have
    \[ 
        \#\left\{\overline{n}\in\mathcal{O}_K/\mathfrak{p}^2 \vcentcolon n\in\mathbb{Z}\right\} = p^2.
    \]
    Consequently,
    \[  
        \#Z_X =\left(\prod_{p\in\mathcal{I}}p^2\right)\left(\prod_{\mathfrak{p}\in\mathcal{P}_X}p^2\right).
    \]
    It follows that
    \[
        \frac{\#N_X}{\#Z_X} = \frac{1}{\displaystyle\prod_{p\in\mathcal{I}}p^2} \prod_{\mathfrak{p}\in\mathcal{P}_X} \left(1-\frac{\rho(\mathfrak{p}^2)}{p^2}\right).
    \]
    
    For finitely many primes $\mathfrak{p}$, we have $\rho(\mathfrak{p}^2)\leq p^2$. For the other primes $\mathfrak{p}$, we have $\rho(\mathfrak{p}^2)\leq2$. Moreover, the corresponding rational primes $p$ are distinct by Proposition~\ref{prop: conjugatenotdividing}. Hence
    \[
        \sum_{\mathfrak{p} \in \mathcal{P}_X} \frac{\rho(\mathfrak{p}^2)}{p^2}
    \]
    converges, and therefore the infinite product
    \[
        \prod_{\mathfrak{p} \in \mathcal{P}_X}\left(1-\frac{\rho(\mathfrak{p}^2)}{p^2}\right)
    \]
    converges to a positive real number. Thus,
    \[
        \lim_{X\to\infty}\frac{\#N_X}{\#Z_X}>0.
    \]
    This proves the proposition.
\end{proof}

\begin{proof}[Proof of Theorem~\ref{thm: main}]
    By Proposition~\ref{prop: sieve}, for each $p\in\mathcal{I}$, there exists an integer $r_p$ such that $h(r_p)\not\equiv 0\pmod{p^2}$. Moreover, there is a positive proportion of integers $n$ satisfying both
    \[
        n\equiv r_p\pmod{p^2\mathcal{O}_K} \qquad \text{for all }p\in\mathcal{I},
    \]
    and
    \[
        \bigl(N_{L/K}(n-\alpha)\bigr) = (m(n))
    \]
    is squarefree away from $\mathcal{I}_K$.
    
    Fix such an integer $n$, and consider the prime ideal factorization of $(m(n))$. Let $\mathfrak{p}$ be a prime ideal appearing in this factorization, and let $p$ be the rational prime below $\mathfrak{p}$.
    
    Suppose first that $p\notin\mathcal{I}$. Since $(m(n))$ is squarefree away from $\mathcal{I}_K$, we have
    \[
        \mathfrak{p}^2\nmid (m(n)).
    \]
    Furthermore, by Proposition~\ref{prop: residue_degree} and Proposition~\ref{prop: conjugatenotdividing}, we have
    \[
        f_{\mathfrak{p}/p}=1
    \]
    and no Galois conjugate of $\mathfrak{p}$ divides $(m(n))$. Therefore, when taking the norm to $\mathbb{Q}$, the prime $p$ occurs with exponent at most one in $N_{K/\mathbb{Q}}((m(n)))= (h(n))$. Hence,
    \[
        p^2\nmid h(n).
    \]
    
    Now suppose that $p\in\mathcal{I}$. Since $n\equiv r_p\pmod{p^2\mathcal{O}_K}$
    and $h(r_p)\not\equiv0\pmod{p^2}$, we have
    \[
        h(n)\equiv h(r_p)\not\equiv0\pmod{p^2}.
    \]
    Thus,
    \[
        p^2\nmid h(n).
    \]
    
    Consequently, $h(n)$ is squarefree for a positive proportion of integers $n$.
\end{proof}

\begin{coro}
    \label{coro:quartics}
    If $h(x) \in \ZZ[x]$ is an irreducible monic polynomial whose Galois group is neither  $S_4$ nor $A_4$ and $\gcd\{h(n) \vcentcolon n\in \ZZ\}$ is squarefree, then $h(n)$ is squarefree for a positive proportion of integers $n$.
\end{coro}

\section{An Application to Unit Shapes}

We now turn to a family of cubic fields known as \emph{exceptional cubic fields}. Our goal is to show that this family contains an infinite subfamily whose unit shapes converge to the hexagonal lattice. To achieve this, we require an explicit system of fundamental units for each field. Conjecture~\ref{conj: enola} predicts the form of such a system, and Theorem~\ref{thm: main} enables us to verify the conjecture for infinitely many exceptional cubic fields. Consequently, we obtain an infinite family of exceptional cubic fields whose unit shapes converge to the hexagonal lattice.



\subsection{Unit Shapes of Number Fields}

For any number field $L$, let $\sigma_1,\dots,\sigma_r$ be the real embeddings of $L$ and $\tau_1, \overline{\tau_1},\dots, \tau_s, \overline{\tau_s}$ be the pairs of complex conjugate embeddings of $L$. Let $E_L$ be the group of units of $\cO_K$ modulo torsion. There is a \emph{logarithmic embedding}
\[
    \Log \colon E_L		  	\longrightarrow \RR^{r + s}	
\]
\[
u		\mapsto (\log\lvert \sigma_1(u) \rvert, \hdots, \log\lvert\sigma_r(u)\rvert, 2\log \lvert \tau_1(u) \rvert,\hdots,2\log \lvert \tau_s(u) \rvert).
\]
Restricting the standard quadratic form on $\RR^{r+s}$ to $E_L$ equips $E_L$ with the structure of a lattice of rank $r + s - 1$. We say $E_L$ is the \emph{unit lattice} of $L$. 

We define the \emph{unit shape} of $L$ to be the equivalence class of this lattice under scaling, reflection, and rotation.  Hence for each such number field $L$ we obtain a point in the moduli space $\mathcal{S}_{r+s-1} \coloneqq \GL_{r+ s-1}(\ZZ)\backslash \GL_{r+ s-1}(\RR) / \GO_{r+ s-1}(\RR)$ of unimodular rank $r + s - 1$ lattices up to homothety and reflection. We call this point the \emph{unit shape} of $L$.


\subsection{Exceptional Cubic Fields}
 A unit $u$ of a number field is said to be \emph{exceptional} if $u-1$ is a unit too. A number field that possesses an exceptional unit is called an \emph{exceptional number field}.

The only exceptional non-Galois real cubic fields are the following infinite family, generated by the exceptional units with minimal polynomials (\cite{nagell}):
\begin{equation}
    \label{eqn:exceptional_cubic}
    f_t(x) \coloneqq x^3 + (t-1)x^2 - tx -1, \qquad t \in \ZZ, t \geq 3.
\end{equation}


\begin{conjecture}[\cite{ennola}]
    \label{conj: enola}
    If $\alpha_t$ is a root of the polynomial~\ref{eqn:exceptional_cubic}, then $\{\alpha_t, \alpha_t -1\}$ is a fundamental system of units for the maximal order of $\QQ(\alpha_t)$.
\end{conjecture}

We can apply Theorem~\ref{thm: main} to show that Conjecture~\ref{conj: enola} is true for a positive proportion of values of $t$ as follows. 

For the remainder of this section, we write \(f(t) \sim g(t)\) for two functions \(f\) and \(g\) if
    \[
        \lim_{t\to\infty}\frac{f(t)}{g(t)}=1.
    \]

\begin{theorem}
    \label{thm: ennola}
    There is a positive proportion of integers $t \geq 3$ for which $\{\alpha_t, \alpha_t - 1\}$ is a fundamental system of units for the maximal order of $\QQ(\alpha_t)$.
\end{theorem}

\begin{lemma}
\label{lem: root-asymptotics}
    If $\alpha_{1,t} < \alpha_{2,t} < \alpha_{3,t}$ are the roots of $f_t(x)$, then
    \[
       \alpha_{1,t} \sim -t + \frac{1}{t^2}, \qquad \alpha_{2,t} \sim -\frac{1}{t}, \qquad \alpha_{3,t} \sim 1 + \frac{1}{t}.
    \]
\end{lemma}
\begin{proof}
    Note that $f_t(x) + 1 = x(x + t)(x-1)$. Therefore, we can use the roots of this polynomial to approximate the roots of $f_t(x)$. Let $\epsilon > 0$ be an arbitrary small real number depending on $t$. 
    \begin{enumerate}
    \item Root $x = -t$.
            \begin{equation*}
                f_t(-t + \epsilon) = (-t+\epsilon)(\epsilon)(-t+\epsilon -1) -1 =t^2\epsilon - 2 t \epsilon^2 + t\epsilon + \epsilon^3 - \epsilon^2 -1
            \end{equation*}
            If we let $\epsilon = \frac{1}{t^2}$, then
            \begin{equation*}
                f_t(-t + \epsilon) = - 2 t \epsilon^2 + t\epsilon + \epsilon^3 - \epsilon^2 \to 0 \text{ as } t \to \infty.
            \end{equation*}
    
        
    \item Root $x = 0$.
    \begin{equation*}
        f_t(\epsilon) = \epsilon(\epsilon + t)(\epsilon -1) -1 = t\epsilon^2 - t\epsilon + \epsilon^3 - \epsilon^2 -1
    \end{equation*}
    If we let $\epsilon = - \frac{1}{t}$, then
    \begin{equation*}
        f_t(\epsilon) =  t\epsilon^2 + \epsilon^3 - \epsilon^2 \to 0 \text{ as } t \to \infty.
    \end{equation*}
    
    \item Root $x = 1$.
    \begin{equation*}
        f_t(1 +\epsilon) = (1 + \epsilon)(1+ \epsilon + t)(\epsilon) -1 = t\epsilon^2 + t\epsilon + \epsilon^3 + 2\epsilon^2 + \epsilon -1.
    \end{equation*}
    If we let $\epsilon = \frac{1}{t}$, then
    \begin{equation*}
        f_t(1+ \epsilon) = t\epsilon^2 + \epsilon^3 + 2\epsilon^2 + \epsilon \to 0 \text{ as } t \to \infty.
    \end{equation*}
    \end{enumerate}
    Finally, by the Intermediate Value Theorem, two of the roots of $f_t(x)$ lie in the intervals $(-1,0)$ and $(0,2)$, respectively, for all $t\geq 3$. Now fix a real number $b > 0$. For sufficiently large $t$,
    \[
        f_t\left( -t + \frac{1}{t^2} - b\right) < 0  \text{ and }  f_t\left( -t + \frac{1}{t^2} + b\right) > 0.
    \]
    By the Intermediate Value Theorem
    \[
        -t + \frac{1}{t^2} - b < \alpha_{1,t} < -t + \frac{1}{t^2} + b.
    \]
    Consequently
    \[
        \alpha_{1,t} \sim -t + \frac{1}{t^2}.
    \]
    A similar reasoning shows that
    \[
         \alpha_{2,t} \sim -\frac{1}{t} \qquad \text{ and } \qquad \alpha_{3,t} \sim 1 + \frac{1}{t}.
    \]
\end{proof}

\begin{proof}[Proof of Theorem~\ref{thm: ennola}]
        Let $\alpha\coloneqq \alpha_{1,t}$. Observe that
        \[
            \alpha-1 \sim -t+\frac{1}{t^2}-1 \sim -t-1.
        \]
    
    The regulator $R$ of the of units $\{\alpha,\alpha-1\}$ is
        \[
            R = \left|\det
            \begin{bmatrix}
                \log|\alpha| & \log|\alpha_2| \\
                \log|\alpha-1| & \log|\alpha_2-1|
            \end{bmatrix}
            \right|
            \sim
            \det
            \begin{bmatrix}
                \log t & -\log t \\
                \log t & 0
            \end{bmatrix}
            =
            (\log t)^2.
        \]
    
    Next, observe that the discriminant of $f_t(x)$ is
        \[
            \Delta(f_t) = t^4 + 6t^3 + 7t^2 - 6t - 31 \sim t^4.
        \]
    
    Viewing $\Delta(f_t)$ as a polynomial in $t$, we see that it satisfies the hypotheses of Corollary~\ref{coro:quartics}. Hence, there is a positive proportion of integers $t$ for which $\Delta(f_t)$ is squarefree. In particular, for such values of $t$,
        \[
            \mathcal{O}_{\mathbb{Q}(\alpha)}  = \mathbb{Z}[\alpha].
        \]
    
    Moreover,
        \[
            \frac{16R}{\log^2(\Delta(f_t)/4)} \sim \frac{16(\log t)^2}{\log^2(t^4/4)} \sim 1.
        \]
    Therefore, Cusick's lower bound for regulators~\cite{cusick-1} implies that $\{\alpha,\alpha-1\}$ is a fundamental system of units of $\mathcal{O}_{\mathbb{Q}(\alpha)}$ for a positive proportion of integers $t$.
\end{proof}


\subsection{Unit Shapes of Exceptional Cubic Fields}

Consider the polynomials $f_t(x)$ for which $\operatorname{Disc}(f_t)$ is squarefree, and let
    \[
        L_t\coloneqq\mathbb{Q}(\alpha_{1,t}).
    \]
By Theorem~\ref{thm: ennola}, the pair $\{\alpha_{1,t},\alpha_{1,t}-1\}$ forms a fundamental system of units of $L_t$. Consequently,
    \[
        \{\Log(\alpha_{1,t}),\Log(\alpha_{1,t}-1)\}
    \]
is a basis of the unit lattice of $L_t$, which is contained in the diagonal hyperplane of $\RR^3$ cut out by $x + y + z =0$. To determine the shape of this lattice, we study the Gram matrix of the vectors $\Log(\alpha_{1,t})$ and $\Log(\alpha_{1,t}-1)$. 

\begin{theorem}
\label{thm: shapes-convergence}
    The unit shapes of the fields $L_t$ converge to the hexagonal lattice.
\end{theorem}
\begin{proof}
     Write $\alpha\coloneqq\alpha_{1,t}, \alpha_2 \coloneqq \alpha_{2,t}$, and $\alpha_3 \coloneqq  \alpha_{3,t}$. Using the asymptotics of Lemma~\ref{lem: root-asymptotics}, we obtain the following. 
    \begin{align*}
           \cos \theta  
           &= \frac{\log|\alpha|\log|\alpha -1| + \log|\alpha_2|\log|\alpha_2-1|+ \log|\alpha_3|\log|\alpha_3-1|}{\sqrt{\log^2|\alpha| + \log^2|\alpha_2| + \log^2|\alpha_3|}\sqrt{\log^2|\alpha-1| + \log^2|\alpha_2-1| + \log^2|\alpha_3-1|} } \\
           &\sim \frac{(\log t)^2}{(\sqrt{2}\log t)(\sqrt{2}\log t)}\\
           & = \frac{1}{2}.       
    \end{align*}

Additionally
    \[
        \|\Log(\alpha)\|^2 = \log^2|\alpha| + \log^2|\alpha_2| + \log^2|\alpha_3| \sim 2(\log t)^2,
    \]
and similarly
    \[
        \|\Log(\alpha-1)\|^2 \sim 2(\log t)^2.
    \]
Thus, we have
    \[
        \frac{\|\Log(\alpha)\|}{\|\Log(\alpha-1)\|} \sim 1.
    \]
After normalizing by a common scaling factor, the Gram matrix of the basis $\{\Log(\alpha),\Log(\alpha-1)\}$ converges to
    \[
        \begin{bmatrix}
        1 & \frac12\\
        \frac12 & 1
        \end{bmatrix},
    \]
which is the Gram matrix of the hexagonal lattice. Therefore, the shapes of the unit lattices of the fields $L_t$ converge to the hexagonal lattice.   
\end{proof}







\begingroup
\hypersetup{
  urlcolor=mybiburlcolor
}
\bibliography{bib.bib}{} \bibliographystyle{my-amsalpha}
\endgroup
\end{document}
